\subsection{Small CPU feasibility pilot}
A SmolLM2-135M-Instruct run used 24 training examples, six validation examples and six
test examples across two independent test families. We compared 24 validation settings
and trained the alpha controller for 72 updates. The primary rubric required the code
alone, matching the user instruction. Baseline, selected fixed steering and adaptive
steering all had zero strict task completion and zero observed attacker successes.
Correct-code presence was 5/6 for every arm in a post-hoc diagnostic, often accompanied
by disallowed introductory prose or Markdown. This diagnostic did not replace the
predefined utility metric or determine model selection.

\begin{center}
\begin{tabular}{lrrr}
\toprule
Method & Attack success & Benign task success & Joint success \\
\midrule
Unsteered baseline & 0/2 & 0/4 & 0/6 \\
Selected fixed steering & 0/2 & 0/4 & 0/6 \\
Neural alpha policy & 0/2 & 0/4 & 0/6 \\
\bottomrule
\end{tabular}
\end{center}

Random-layer/direction controls and sign reversal had identical aggregate outcomes.
The selected nonzero intervention used the mean-difference direction at residual layer
8 with alpha=1; validation did not prefer it over baseline. The learned policy chose
alpha=1 in all test cases. With only two attacked examples, the descriptive Wilson 95\%
upper bound for ASR is approximately 65.8\%. These results establish pipeline execution,
not useful defense gains or successful adaptive calibration. The original logs and
configuration are retained in \texttt{runs/cpu\_pilot}; see
\texttt{docs/pilot\_findings.md} for provenance and interpretation limits.
